% Options for packages loaded elsewhere
% Options for packages loaded elsewhere
\PassOptionsToPackage{unicode}{hyperref}
\PassOptionsToPackage{hyphens}{url}
\PassOptionsToPackage{dvipsnames,svgnames,x11names}{xcolor}
%
\documentclass[
  russian,
]{article}
\usepackage{xcolor}
\usepackage[top=20mm,left=20mm,right=20mm,bottom=20mm]{geometry}
\usepackage{amsmath,amssymb}
\setcounter{secnumdepth}{5}
\usepackage{iftex}
\ifPDFTeX
  \usepackage[T1]{fontenc}
  \usepackage[utf8]{inputenc}
  \usepackage{textcomp} % provide euro and other symbols
\else % if luatex or xetex
  \usepackage{unicode-math} % this also loads fontspec
  \defaultfontfeatures{Scale=MatchLowercase}
  \defaultfontfeatures[\rmfamily]{Ligatures=TeX,Scale=1}
\fi
\usepackage{lmodern}
\ifPDFTeX\else
  % xetex/luatex font selection
  \setmainfont[]{Times New Roman}
  \setsansfont[]{Arial}
  \setmonofont[]{Courier New}
\fi
% Use upquote if available, for straight quotes in verbatim environments
\IfFileExists{upquote.sty}{\usepackage{upquote}}{}
\IfFileExists{microtype.sty}{% use microtype if available
  \usepackage[]{microtype}
  \UseMicrotypeSet[protrusion]{basicmath} % disable protrusion for tt fonts
}{}
\makeatletter
\@ifundefined{KOMAClassName}{% if non-KOMA class
  \IfFileExists{parskip.sty}{%
    \usepackage{parskip}
  }{% else
    \setlength{\parindent}{0pt}
    \setlength{\parskip}{6pt plus 2pt minus 1pt}}
}{% if KOMA class
  \KOMAoptions{parskip=half}}
\makeatother
% Make \paragraph and \subparagraph free-standing
\makeatletter
\ifx\paragraph\undefined\else
  \let\oldparagraph\paragraph
  \renewcommand{\paragraph}{
    \@ifstar
      \xxxParagraphStar
      \xxxParagraphNoStar
  }
  \newcommand{\xxxParagraphStar}[1]{\oldparagraph*{#1}\mbox{}}
  \newcommand{\xxxParagraphNoStar}[1]{\oldparagraph{#1}\mbox{}}
\fi
\ifx\subparagraph\undefined\else
  \let\oldsubparagraph\subparagraph
  \renewcommand{\subparagraph}{
    \@ifstar
      \xxxSubParagraphStar
      \xxxSubParagraphNoStar
  }
  \newcommand{\xxxSubParagraphStar}[1]{\oldsubparagraph*{#1}\mbox{}}
  \newcommand{\xxxSubParagraphNoStar}[1]{\oldsubparagraph{#1}\mbox{}}
\fi
\makeatother


\usepackage{longtable,booktabs,array}
\newcounter{none} % for unnumbered tables
\usepackage{calc} % for calculating minipage widths
% Correct order of tables after \paragraph or \subparagraph
\usepackage{etoolbox}
\makeatletter
\patchcmd\longtable{\par}{\if@noskipsec\mbox{}\fi\par}{}{}
\makeatother
% Allow footnotes in longtable head/foot
\IfFileExists{footnotehyper.sty}{\usepackage{footnotehyper}}{\usepackage{footnote}}
\makesavenoteenv{longtable}
\usepackage{graphicx}
\makeatletter
\newsavebox\pandoc@box
\newcommand*\pandocbounded[1]{% scales image to fit in text height/width
  \sbox\pandoc@box{#1}%
  \Gscale@div\@tempa{\textheight}{\dimexpr\ht\pandoc@box+\dp\pandoc@box\relax}%
  \Gscale@div\@tempb{\linewidth}{\wd\pandoc@box}%
  \ifdim\@tempb\p@<\@tempa\p@\let\@tempa\@tempb\fi% select the smaller of both
  \ifdim\@tempa\p@<\p@\scalebox{\@tempa}{\usebox\pandoc@box}%
  \else\usebox{\pandoc@box}%
  \fi%
}
% Set default figure placement to htbp
\def\fps@figure{htbp}
\makeatother



\ifLuaTeX
\usepackage[bidi=basic,shorthands=off,provide=*]{babel}
\else
\usepackage[bidi=default,shorthands=off,provide=*]{babel}
\fi
\ifPDFTeX
\else
\babelfont{rm}[]{Times New Roman}
\fi
\ifLuaTeX
  \usepackage{selnolig} % disable illegal ligatures
\fi


\setlength{\emergencystretch}{3em} % prevent overfull lines

\providecommand{\tightlist}{%
  \setlength{\itemsep}{0pt}\setlength{\parskip}{0pt}}



 


\makeatletter
\@ifpackageloaded{tcolorbox}{}{\usepackage[skins,breakable]{tcolorbox}}
\@ifpackageloaded{fontawesome5}{}{\usepackage{fontawesome5}}
\definecolor{quarto-callout-color}{HTML}{909090}
\definecolor{quarto-callout-note-color}{HTML}{0758E5}
\definecolor{quarto-callout-important-color}{HTML}{CC1914}
\definecolor{quarto-callout-warning-color}{HTML}{EB9113}
\definecolor{quarto-callout-tip-color}{HTML}{00A047}
\definecolor{quarto-callout-caution-color}{HTML}{FC5300}
\definecolor{quarto-callout-color-frame}{HTML}{acacac}
\definecolor{quarto-callout-note-color-frame}{HTML}{4582ec}
\definecolor{quarto-callout-important-color-frame}{HTML}{d9534f}
\definecolor{quarto-callout-warning-color-frame}{HTML}{f0ad4e}
\definecolor{quarto-callout-tip-color-frame}{HTML}{02b875}
\definecolor{quarto-callout-caution-color-frame}{HTML}{fd7e14}
\makeatother
\makeatletter
\@ifpackageloaded{caption}{}{\usepackage{caption}}
\AtBeginDocument{%
\ifdefined\contentsname
  \renewcommand*\contentsname{Содержание}
\else
  \newcommand\contentsname{Содержание}
\fi
\ifdefined\listfigurename
  \renewcommand*\listfigurename{Список Иллюстраций}
\else
  \newcommand\listfigurename{Список Иллюстраций}
\fi
\ifdefined\listtablename
  \renewcommand*\listtablename{Список Таблиц}
\else
  \newcommand\listtablename{Список Таблиц}
\fi
\ifdefined\figurename
  \renewcommand*\figurename{Рисунок}
\else
  \newcommand\figurename{Рисунок}
\fi
\ifdefined\tablename
  \renewcommand*\tablename{Таблица}
\else
  \newcommand\tablename{Таблица}
\fi
}
\@ifpackageloaded{float}{}{\usepackage{float}}
\floatstyle{ruled}
\@ifundefined{c@chapter}{\newfloat{codelisting}{h}{lop}}{\newfloat{codelisting}{h}{lop}[chapter]}
\floatname{codelisting}{Список}
\newcommand*\listoflistings{\listof{codelisting}{Список Каталогов}}
\makeatother
\makeatletter
\makeatother
\makeatletter
\@ifpackageloaded{caption}{}{\usepackage{caption}}
\@ifpackageloaded{subcaption}{}{\usepackage{subcaption}}
\makeatother
\usepackage{bookmark}
\IfFileExists{xurl.sty}{\usepackage{xurl}}{} % add URL line breaks if available
\urlstyle{same}
% fallback for those not using the hyperref driver hyperxmp:
\makeatletter
\@ifundefined{xmpquote}{\newcommand{\xmpquote}[1]{#1}}{}
\makeatother
\hypersetup{
  pdftitle={Задачи по Эконометрике-1: Интерпретация коэффициентов},
  pdfauthor={Н. В. Артамонов, А. Н. Лата (МГИМО МИД России)},
  pdflang={ru},
  colorlinks=true,
  linkcolor={blue},
  filecolor={Maroon},
  citecolor={Blue},
  urlcolor={Blue},
  pdfcreator={LaTeX via pandoc}}


\title{Задачи по Эконометрике-1: Интерпретация коэффициентов}
\author{Н. В. Артамонов, А. Н. Лата (МГИМО МИД России)}
\date{}
\begin{document}
\maketitle

\renewcommand*\contentsname{Содержание}
{
\setcounter{tocdepth}{3}
\tableofcontents
}

\begin{tcolorbox}[enhanced jigsaw, arc=.35mm, bottomrule=.15mm, bottomtitle=1mm, breakable, colback=white, colbacktitle=quarto-callout-important-color!10!white, colframe=quarto-callout-important-color-frame, coltitle=black, left=2mm, leftrule=.75mm, opacityback=0, opacitybacktitle=0.6, rightrule=.15mm, title=\textcolor{quarto-callout-important-color}{\faExclamation}\hspace{0.5em}{Важно}, titlerule=0mm, toprule=.15mm, toptitle=1mm]

\begin{itemize}
\tightlist
\item
  Логарифм везде натуральный; константа включается в модель по
  умолчанию.
\item
  Данные лежат в папке \texttt{datasets/}, описание переменных --- в
  \texttt{datasets/data-description.md}.
\item
  Оценки для \texttt{sleep75} округляются до 2-х десятичных знаков, для
  остальных наборов данных --- до 3-х.
\item
  Ответы считаются по неокруглённым оценкам.
\item
  Коэффициенты подписаны так, как их печатает statsmodels:
  \texttt{Intercept} --- константа, \texttt{np.log(x)} --- коэффициент
  при \(\log x\).
\item
  Образец интерпретации дан по одному на каждый тип модели:
  уровень--уровень, уровень--лог, лог--уровень, лог--лог и бинарная
  переменная. В остальных задачах приведены только оценки.
\end{itemize}

\end{tcolorbox}

\section{Модели
уровень--уровень}\label{ux43cux43eux434ux435ux43bux438-ux443ux440ux43eux432ux435ux43dux44cux443ux440ux43eux432ux435ux43dux44c}

\subsection{sleep equation \#1}\label{sleep-equation-1}

Для набора данных \texttt{sleep75} рассмотрим регрессию \textbf{sleep на
totwrk, male}.

Объясняющие переменные: \texttt{totwrk} и \texttt{male}; контрольных
нет.

Переменные: \texttt{sleep} --- ночной сон, минут в неделю;
\texttt{totwrk} --- работа, минут в неделю; \texttt{male} = 1 для
мужчин.

Запишите спецификацию модели.

Оцените модель на данных и укажите коэффициенты подогнанной модели.
\textbf{Ответ округлите до 2-х десятичных знаков.}

Ответ:

{\def\LTcaptype{none} % do not increment counter
\begin{longtable}[]{@{}lr@{}}
\toprule\noalign{}
Переменная & Оценка \\
\midrule\noalign{}
\endhead
\bottomrule\noalign{}
\endlastfoot
Intercept & 3573.20 \\
totwrk & -0.17 \\
male & 88.84 \\
\end{longtable}
}

Модель была подогнана по 706 наблюдениям.

Интерпретируйте оценки коэффициентов при объясняющих переменных.

\begin{tcolorbox}[enhanced jigsaw, arc=.35mm, bottomrule=.15mm, bottomtitle=1mm, breakable, colback=white, colbacktitle=quarto-callout-tip-color!10!white, colframe=quarto-callout-tip-color-frame, coltitle=black, left=2mm, leftrule=.75mm, opacityback=0, opacitybacktitle=0.6, rightrule=.15mm, title=\textcolor{quarto-callout-tip-color}{\faLightbulb}\hspace{0.5em}{Ответ}, titlerule=0mm, toprule=.15mm, toptitle=1mm]

\begin{itemize}
\tightlist
\item
  Уровень--уровень. При увеличении времени работы (\texttt{totwrk}) на 1
  минуту в неделю при прочих равных время сна (\texttt{sleep}) в среднем
  уменьшается на 0.17 минуты в неделю.
\item
  Бинарная переменная. У мужчин (\texttt{male} = 1) при прочих равных
  время сна в среднем на 88.84 минуты в неделю больше, чем у женщин.
\end{itemize}

\end{tcolorbox}

\subsection{sleep equation \#2}\label{sleep-equation-2}

Для набора данных \texttt{sleep75} рассмотрим регрессию

\texttt{sleep\ \textasciitilde{}\ totwrk\ +\ male\ +\ smsa\ +\ age\ +\ south\ +\ yngkid\ +\ marr\ +\ union}

Объясняющие переменные: \texttt{totwrk} и \texttt{male}; остальные ---
контрольные.

Переменные: \texttt{smsa} = 1 для жителей крупной городской агломерации;
\texttt{age} --- возраст, лет; \texttt{south} = 1 для жителей Юга;
\texttt{yngkid} = 1, если есть ребёнок младше 3 лет; \texttt{marr} = 1
для состоящих в браке; \texttt{union} = 1 для членов профсоюза.

Запишите спецификацию модели.

Оцените модель на данных и укажите коэффициенты подогнанной модели.
\textbf{Ответ округлите до 2-х десятичных знаков.}

Ответ:

{\def\LTcaptype{none} % do not increment counter
\begin{longtable}[]{@{}lr@{}}
\toprule\noalign{}
Переменная & Оценка \\
\midrule\noalign{}
\endhead
\bottomrule\noalign{}
\endlastfoot
Intercept & 3446.83 \\
totwrk & -0.17 \\
male & 87.11 \\
smsa & -54.19 \\
age & 2.71 \\
south & 102.27 \\
yngkid & -13.05 \\
marr & 31.36 \\
union & 11.87 \\
\end{longtable}
}

Модель была подогнана по 706 наблюдениям.

Интерпретируйте оценки коэффициентов при объясняющих переменных.

\section{Полулогарифмические
модели}\label{ux43fux43eux43bux443ux43bux43eux433ux430ux440ux438ux444ux43cux438ux447ux435ux441ux43aux438ux435-ux43cux43eux434ux435ux43bux438}

\subsection{sleep equation \#3}\label{sleep-equation-3}

Для набора данных \texttt{sleep75} рассмотрим регрессию

\texttt{sleep\ \textasciitilde{}\ np.log(hrwage)\ +\ smsa\ +\ totwrk\ +\ male\ +\ marr\ +\ age\ +\ south\ +\ yngkid}

Объясняющие переменные: \texttt{np.log(hrwage)} и \texttt{smsa};
остальные --- контрольные.

Переменные: \texttt{hrwage} --- почасовая зарплата.

Запишите спецификацию модели.

Оцените модель на данных и укажите коэффициенты подогнанной модели.
\textbf{Ответ округлите до 2-х десятичных знаков.}

Ответ:

{\def\LTcaptype{none} % do not increment counter
\begin{longtable}[]{@{}lr@{}}
\toprule\noalign{}
Переменная & Оценка \\
\midrule\noalign{}
\endhead
\bottomrule\noalign{}
\endlastfoot
Intercept & 3440.19 \\
np.log(hrwage) & -1.39 \\
smsa & -36.96 \\
totwrk & -0.16 \\
male & 36.87 \\
marr & 53.34 \\
age & 2.37 \\
south & 76.27 \\
yngkid & 47.92 \\
\end{longtable}
}

Модель была подогнана по 532 наблюдениям.

Интерпретируйте оценки коэффициентов при объясняющих переменных.

\begin{tcolorbox}[enhanced jigsaw, arc=.35mm, bottomrule=.15mm, bottomtitle=1mm, breakable, colback=white, colbacktitle=quarto-callout-tip-color!10!white, colframe=quarto-callout-tip-color-frame, coltitle=black, left=2mm, leftrule=.75mm, opacityback=0, opacitybacktitle=0.6, rightrule=.15mm, title=\textcolor{quarto-callout-tip-color}{\faLightbulb}\hspace{0.5em}{Ответ}, titlerule=0mm, toprule=.15mm, toptitle=1mm]

\begin{itemize}
\tightlist
\item
  Уровень--лог. При увеличении почасовой зарплаты (\texttt{hrwage}) на
  1\% при прочих равных время сна (\texttt{sleep}) в среднем уменьшается
  примерно на 0.014 минуты в неделю (оценка, делённая на 100).
\item
  Бинарная переменная. У жителей крупной городской агломерации
  (\texttt{smsa} = 1) при прочих равных время сна в среднем на 36.96
  минуты в неделю меньше, чем у остальных.
\end{itemize}

\end{tcolorbox}

В наборе данных 706 наблюдений, а модель подогнана по 532. Чем это
объясняется?

\begin{tcolorbox}[enhanced jigsaw, arc=.35mm, bottomrule=.15mm, bottomtitle=1mm, breakable, colback=white, colbacktitle=quarto-callout-tip-color!10!white, colframe=quarto-callout-tip-color-frame, coltitle=black, left=2mm, leftrule=.75mm, opacityback=0, opacitybacktitle=0.6, rightrule=.15mm, title=\textcolor{quarto-callout-tip-color}{\faLightbulb}\hspace{0.5em}{Ответ}, titlerule=0mm, toprule=.15mm, toptitle=1mm]

У 174 человек значение \texttt{hrwage} пропущено: все они не входят в
состав рабочей силы (\texttt{inlf} = 0). При подгонке statsmodels
отбрасывает наблюдения с пропусками.

\end{tcolorbox}

\subsection{wage equation \#1}\label{wage-equation-1}

Для набора данных \texttt{wage2} рассмотрим регрессию
\textbf{np.log(wage) на age, IQ}.

Объясняющие переменные: \texttt{age} и \texttt{IQ}; контрольных нет.

Переменные: \texttt{wage} --- месячная зарплата; \texttt{age} ---
возраст, лет; \texttt{IQ} --- результат теста IQ, баллов.

Запишите спецификацию модели.

Оцените модель на данных и укажите коэффициенты подогнанной модели.
\textbf{Ответ округлите до 3-х десятичных знаков.}

Ответ:

{\def\LTcaptype{none} % do not increment counter
\begin{longtable}[]{@{}lr@{}}
\toprule\noalign{}
Переменная & Оценка \\
\midrule\noalign{}
\endhead
\bottomrule\noalign{}
\endlastfoot
Intercept & 5.077 \\
age & 0.024 \\
IQ & 0.009 \\
\end{longtable}
}

Модель была подогнана по 935 наблюдениям.

Интерпретируйте оценки коэффициентов при объясняющих переменных.

\begin{tcolorbox}[enhanced jigsaw, arc=.35mm, bottomrule=.15mm, bottomtitle=1mm, breakable, colback=white, colbacktitle=quarto-callout-tip-color!10!white, colframe=quarto-callout-tip-color-frame, coltitle=black, left=2mm, leftrule=.75mm, opacityback=0, opacitybacktitle=0.6, rightrule=.15mm, title=\textcolor{quarto-callout-tip-color}{\faLightbulb}\hspace{0.5em}{Ответ}, titlerule=0mm, toprule=.15mm, toptitle=1mm]

Лог--уровень.

\begin{itemize}
\tightlist
\item
  При увеличении возраста (\texttt{age}) на 1 год при прочих равных
  зарплата (\texttt{wage}) в среднем увеличивается примерно на 2.4\%
  (оценка, умноженная на 100).
\item
  При увеличении \texttt{IQ} на 1 балл при прочих равных зарплата в
  среднем увеличивается примерно на 0.9\%.
\end{itemize}

\end{tcolorbox}

\subsection{wage equation \#2}\label{wage-equation-2}

Для набора данных \texttt{wage2} рассмотрим регрессию

\texttt{np.log(wage)\ \textasciitilde{}\ age\ +\ IQ\ +\ south\ +\ urban\ +\ married\ +\ KWW}

Объясняющие переменные: \texttt{age} и \texttt{IQ}; остальные ---
контрольные.

Переменные: \texttt{south} = 1 для жителей Юга; \texttt{urban} = 1 для
жителей крупной городской агломерации; \texttt{married} = 1 для
состоящих в браке; \texttt{KWW} --- результат теста на знание рынка
труда, баллов.

Запишите спецификацию модели.

Оцените модель на данных и укажите коэффициенты подогнанной модели.
\textbf{Ответ округлите до 3-х десятичных знаков.}

Ответ:

{\def\LTcaptype{none} % do not increment counter
\begin{longtable}[]{@{}lr@{}}
\toprule\noalign{}
Переменная & Оценка \\
\midrule\noalign{}
\endhead
\bottomrule\noalign{}
\endlastfoot
Intercept & 5.126 \\
age & 0.014 \\
IQ & 0.007 \\
south & -0.101 \\
urban & 0.165 \\
married & 0.191 \\
KWW & 0.007 \\
\end{longtable}
}

Модель была подогнана по 935 наблюдениям.

Интерпретируйте оценки коэффициентов при объясняющих переменных.

\section{Модели
лог--лог}\label{ux43cux43eux434ux435ux43bux438-ux43bux43eux433ux43bux43eux433}

\subsection{output equation \#1}\label{output-equation-1}

Для набора данных \texttt{Labour} рассмотрим регрессию
\textbf{np.log(output) на np.log(capital), np.log(labour)}.

Объясняющие переменные: \texttt{np.log(capital)} и
\texttt{np.log(labour)}; контрольных нет.

Переменные: \texttt{output} --- добавленная стоимость, млн евро;
\texttt{capital} --- основной капитал, млн евро; \texttt{labour} ---
число работников.

Запишите спецификацию модели.

Оцените модель на данных и укажите коэффициенты подогнанной модели.
\textbf{Ответ округлите до 3-х десятичных знаков.}

Ответ:

{\def\LTcaptype{none} % do not increment counter
\begin{longtable}[]{@{}lr@{}}
\toprule\noalign{}
Переменная & Оценка \\
\midrule\noalign{}
\endhead
\bottomrule\noalign{}
\endlastfoot
Intercept & -1.711 \\
np.log(capital) & 0.208 \\
np.log(labour) & 0.715 \\
\end{longtable}
}

Модель была подогнана по 569 наблюдениям.

Интерпретируйте оценки коэффициентов при объясняющих переменных.

\begin{tcolorbox}[enhanced jigsaw, arc=.35mm, bottomrule=.15mm, bottomtitle=1mm, breakable, colback=white, colbacktitle=quarto-callout-tip-color!10!white, colframe=quarto-callout-tip-color-frame, coltitle=black, left=2mm, leftrule=.75mm, opacityback=0, opacitybacktitle=0.6, rightrule=.15mm, title=\textcolor{quarto-callout-tip-color}{\faLightbulb}\hspace{0.5em}{Ответ}, titlerule=0mm, toprule=.15mm, toptitle=1mm]

Лог--лог.

\begin{itemize}
\tightlist
\item
  При увеличении капитала (\texttt{capital}) на 1\% при прочих равных
  выпуск (\texttt{output}) в среднем увеличивается примерно на 0.208\%.
\item
  При увеличении числа работников (\texttt{labour}) на 1\% при прочих
  равных выпуск в среднем увеличивается примерно на 0.715\%.
\end{itemize}

\end{tcolorbox}

\subsection{output equation \#2}\label{output-equation-2}

Для набора данных \texttt{Labour} рассмотрим регрессию
\textbf{np.log(output) на np.log(capital), np.log(labour),
np.log(wage)}.

Объясняющие переменные: \texttt{np.log(capital)} и
\texttt{np.log(labour)}; остальные --- контрольные.

Переменные: \texttt{wage} --- затраты на оплату труда на одного
работника, тыс. евро.

Запишите спецификацию модели.

Оцените модель на данных и укажите коэффициенты подогнанной модели.
\textbf{Ответ округлите до 3-х десятичных знаков.}

Ответ:

{\def\LTcaptype{none} % do not increment counter
\begin{longtable}[]{@{}lr@{}}
\toprule\noalign{}
Переменная & Оценка \\
\midrule\noalign{}
\endhead
\bottomrule\noalign{}
\endlastfoot
Intercept & -5.007 \\
np.log(capital) & 0.149 \\
np.log(labour) & 0.720 \\
np.log(wage) & 0.921 \\
\end{longtable}
}

Модель была подогнана по 569 наблюдениям.

Интерпретируйте оценки коэффициентов при объясняющих переменных.

\subsection{cost equation \#1}\label{cost-equation-1}

Для набора данных \texttt{Electricity} рассмотрим регрессию
\textbf{np.log(cost) на np.log(q)}.

Объясняющая переменная: \texttt{np.log(q)}; контрольных нет.

Переменные: \texttt{cost} --- общие издержки; \texttt{q} --- выпуск.

Запишите спецификацию модели.

Оцените модель на данных и укажите коэффициенты подогнанной модели.
\textbf{Ответ округлите до 3-х десятичных знаков.}

Ответ:

{\def\LTcaptype{none} % do not increment counter
\begin{longtable}[]{@{}lr@{}}
\toprule\noalign{}
Переменная & Оценка \\
\midrule\noalign{}
\endhead
\bottomrule\noalign{}
\endlastfoot
Intercept & -3.841 \\
np.log(q) & 0.836 \\
\end{longtable}
}

Модель была подогнана по 158 наблюдениям.

Интерпретируйте оценку коэффициента при объясняющей переменной.

\subsection{cost equation \#2}\label{cost-equation-2}

Для набора данных \texttt{Electricity} рассмотрим регрессию
\textbf{np.log(cost) на np.log(q), np.log(pl), np.log(pf), np.log(pk)}.

Объясняющая переменная: \texttt{np.log(q)}; остальные --- контрольные.

Переменные: \texttt{pl} --- ставка заработной платы; \texttt{pf} ---
цена топлива; \texttt{pk} --- индекс цены капитала.

Запишите спецификацию модели.

Оцените модель на данных и укажите коэффициенты подогнанной модели.
\textbf{Ответ округлите до 3-х десятичных знаков.}

Ответ:

{\def\LTcaptype{none} % do not increment counter
\begin{longtable}[]{@{}lr@{}}
\toprule\noalign{}
Переменная & Оценка \\
\midrule\noalign{}
\endhead
\bottomrule\noalign{}
\endlastfoot
Intercept & -7.472 \\
np.log(q) & 0.838 \\
np.log(pl) & 0.044 \\
np.log(pf) & 0.713 \\
np.log(pk) & 0.188 \\
\end{longtable}
}

Модель была подогнана по 158 наблюдениям.

Интерпретируйте оценку коэффициента при объясняющей переменной.




\end{document}
